\section{Geocentrism and the Center of the Universe}

\begin{quotex}
God is an infinite sphere, the centre of which is everywhere, the circumference nowhere.
\flright{\textsc{Hermes Trismegistus}}

God is an intelligible sphere, whose centre is everywhere, and whose circumference is nowhere.
\flright{\textsc{Alain de Lille}}

The whole visible world is only an imperceptible atom in the ample bosom of nature. No idea approaches it. We may enlarge our conceptions beyond all imaginable space; we only produce atoms in comparison with the reality of things. It is an infinite sphere, the center of which is everywhere, the circumference nowhere. In short, it is the greatest sensible mark of the almighty power of God that imagination loses itself in that thought.
\flright{\textsc{Blaise Pascal}}
\end{quotex}


A presumptuous truism of modernity is that science has disproved the geocentric theory of the universe. This allegedly makes man aware of his cosmic insignificance, because of which he is always in a state of angst and feels the urge to act like a character in a Woody Allen film.

Actually, the opposite is the case. This diagram of the Empyrean\footnote{\url{https://gornahoor.net/wp-content/uploads/2011/03/hierarchy.jpg}}, properly understood, shows that the earth is the farthest point from God, which makes it insignificant and godforsaken. Only the subterrestrial hell realms are worse and man's continual descent leads him in that direction. This thought brought the Traditional Man to anxiety; hence, his concern with salvation and ultimate deliverance.

Man, in his descent from his true Self, is thrown into the historico-material process in which he forgets his own being. Modern man can only see the material or historical process, or sometimes both together, which he calls the problem of ``nature vs nurture''. He conceives man as a \emph{tabula rasa}, that is, lacking a nature, karma, dharma. Hence, his thrownness is an injustice, his defects and strengths are not his, but cosmic accidents. Nor can he claim credit for his intelligence, wisdom, courage, fortitude or other qualities. His defects need to be corrected through the intervention of others.

Words\footnote{\url{https://gornahoor.net/?p=3936}}, to the modern mind, are sophisticated animal grunts and have no relation to anything transcendental. Hence, the intellect leads to no ultimate understanding, but rather is a tool of the will to power. To assuage his guilt over his undeserved advantages, he needs to be committed to or engaged in a cause. The absurdity of the cause means nothing, because the world itself is absurd. That is why it is futile to debate such a person; there is no common ground.

Traditional man never understood the Earth as the center of the Universe, to believe otherwise is risible. Modern men, and those who think like them, understand things in material terms. Hence, they are forced to interpret the diagrams of sacred science in material terms, i.e., as primitive forms of astronomy. No, the Empyrean represents a spiritual journey through states of consciousness, not a physical journey like Star Trek. Man descends through denser states to the sublunar region when he ultimately reaches the subtle and gross states represented by water and earth. Here, his dharma intersects with Fate and Fortuna at his conception and birth.

The subtle state, or water element, of the world is chaotic, consciousness is perturbed with pointless desires, confusions, absurdities that the child absorbs. Most men find it impossible to maintain a steady mind, but that is the only way to allow higher influences to come into consciousness.

The gross state, or earth element, is dominated by blind, deterministic forces. Instead of being the field where a man can work out his karmic influences, it instead proves to be a privation, where his undeveloped will is weak and unable to manipulate and control those forces. Only by intense work on developing his True Will can he hope to progress.

It is clear, then, that the Earth is not our true center. The attempts of the neo-pagans to relocate Tradition in nature are thus misguided. Likewise, are the attempts to support tradition on neo-Nietzscheanism. That is merely another form of modernism; the only result is a battle for the control of the Mythos, but as the will to power, not as the fruit of intellectual understanding.

As \textbf{Rene Guenon}\footnote{\url{https://gornahoor.net/?p=3936}} points out, the ``Heart of the Incarnate Word'' is the ``spiritual Sun and Center of the World''. It is our task to find our way back home to our center, and ``know the place for the first time.'' By finding the center of the World, we will find the God whose center is everywhere.


\flrightit{Posted on 2012-04-10 by Cologero}

\begin{center}* * *\end{center}

\begin{footnotesize}
\begin{sffamily}
\texttt{Avery on 2012-04-11 at 00:07 said:}

Here is how I understood your article:

The materialist, modern observer sees Nature as the origin of all things; the Earth is the center of the universe, not in the superficial sense, but in the deeper sense that claims and values lack substance unless if they relate back to material-scientific law. Language is to modern man a tool used by the strong to dominate the weak for purposes of survival and reproduction.

Primitive man did not use language to bring himself ``down to Earth'', but instead to elevate himself facing God. The spiritless, wordless chaos of Nature is a state of alienation from God which one must avoid, and which is as bad as it gets without actually falling into Hell and facing God's displeasure. Earthly things, especially the matters of sex and violence which now dominate modern discourse, are at best part of a system that enables higher seeking, and at worst a distraction from it.

Please correct any errors in my summary.

\hfill

\texttt{Cologero on 2012-04-11 at 07:44 said:}

We use the term ``primitive man'' is a specific way here; I believe you mean to say ``Traditional man''. ``Primitive'' implies an evolutionary development from a lower state to a higher state, in which ``primitive man'' offers a glimpse into our past. Rather, it represents a glimpse into our future.

\hfill

\texttt{Cologero on 2012-04-11 at 07:49 said:}

Here are some homework questions:

1) How does Alain de Lille, a Medieval theologian, come up with the same formulation as the much more ancient Hermetic Tradition?

2) For Christians: why will you never hear a Sunday sermon on the ideas of Alain de Lille?

3) For new rightists: is this what you mean by the Semitic desert god?

\hfill

\texttt{Graham on 2012-04-11 at 16:58 said:}

Answers to homework questions:

1) Through an unbroken intellectual tradition. For the record I have read the same formulation in Nicholas of Cusa.

2) Because that wouldn't be ``pastoral''. Well, I am only half-joking. Because the knowledge has been cast aside. And what was intuitive to a peasant is obscure and embarrassing to a businessman.

\hfill

\texttt{Cologero on 2012-04-12 at 13:31 said:}

It is a strange characteristic of human nature that a man with a certain character flaw is likely to interpret other men in terms of that flaw. For example, a cheat is more likely to regard all men as cheats. In particular, a man who constantly accuses us of using ``straw man'' arguments, uses nothing but them himself.

We have never promoted any particular dogma, creed, faith, or belief. To the contrary, we have pointed out they are necessarily matters of opinion, not of principle. Nevertheless, we do not consider it wrong for a man to take one or the other side of such issues, for good reasons, provided he understands the difference.

It is true that our focus is on the Western traditions, including Hermetism, antique paganism, and Medievalism. The Eastern traditions are well discussed elsewhere, while, oddly enough, our own Western traditions are poorly understood. This can be attributed to a couple of factors. There is a blindness brought on by polemics. This is revealed by the attempt to prefer one specific belief system over another for reasons that have nothing to do with principles. For example, they may object to an alleged cast of mind, attributed typically to some despised historical people. However, as we, following Evola, have pointed out, that is on the horizontal level of subtle manifestation. This sort of objection has nothing to do with principles that transcend space and time.

The more important factor is that it truly is difficult to understand these principles. Unfortunately, in our era where there are no standards, and the good and the bad are mixed indiscriminately, anyone can claim authority on anything. To be clear, we rely on authority ourselves as a sort of intellectual check. So when we focus on the Middle Ages, we do nothing more than develop what the main authorities of Tradition have written about that time.

Specifically, Guenon, Schuon, Coomaraswamy, de Giorgio, and Evola all regard the Middle Ages as a Traditional civilization, the last one experienced in the West. Guenon and Schuon were non-Christians who regard it as traditional. Guenon shows the correspondence of Medieval principles to those of the East. Coomaraswamy, also a non-Christian (although he advised his son to adhere to it), also shows that correspondence in great detail. Furthermore, he recommends several important Western and Christian texts to understand Tradition. De Giorgio was more sympathetic to Christianity, but also to antique paganism. He demonstrates the connection and continuity between the two. Only Evola is expressly anti-Christian, yet he recognizes the Traditional features of the Middle Ages. He also severely criticized neo-paganism, because he understood it is a matter of principle, not labels or group identity.

So, the typical idiots continue to try to post here, believing they can insult, humiliate, or refute us. Now, we may get some facts wrong, or express things unclearly, or even misrepresent Tradition. Those are valid points for discussion. But to send us a list of alleged facts, which are no more than innuendo, claiming to refute thereby the Traditional character of the Middle Ages is a waste of effort. That is not critiquing us; that is a critique on the very notion of Tradition and its founding fathers as we pointed out in the preceding paragraph.

I add this long comment, not for the benefit of one person who most likely will not see it that way, but on the assumption that there may be others who suffer from the same misunderstanding. It is unfortunate to see a man dedicate a lot of time and effort, and have nothing to show for. Thus, we make the effort to always focus on principles. Those who see life in just two dimensions are forced to engage in Strife. Yet, there is something deeper than that.

\hfill

\texttt{Cologero on 2012-04-12 at 17:09 said:}

To demonstrate the confusion of today and the inadequacy of our common ways of categorizing ideas, here is a quote from a neo-pagan web site:

\begin{quotex}\footnotesize

\textbf{Knut Hamsun} and \textbf{John Tolkien} are perhaps the writers who can best be defined as ``Pagan'' of the last 150 years, pagan and profoundly anti-modern.
\end{quotex}


I know nothing about Mr. Hamsun, but I am certain that Mr. Tolkien would be quite surprised at that judgment, except perhaps for the ``anti-modern'' part. I know he opposed the political system espoused on that web site.

Yet it is an interesting point that illustrates, rather than refutes, our own perspective. We will have more to say about it at some future date. Can the case be made that Lord of the Rings is the most important pagan literature of the past century?

\hfill

\texttt{Matt on 2012-04-12 at 20:01 said:}

I can see they placed Hamsun in that regard, since his noted pantheism seems to be an integral part of most neo-pagan thinking.

\hfill

\texttt{Cologero on 2012-04-12 at 22:11 said:}

But they link the two and interpret LOTR in a naturalistic sense. They say the original is in Italian :

\begin{quotex}\footnotesize

When we read their books, we truly feel those living and vivifying descriptions of nature, her forces, her elements. We see that to speak in the same way of thinking, in brief, it speaks the Hyperborean blood itself.
\end{quotex}


Yet, Tolkien himself writes:

\begin{quotex}\footnotesize

The Lord of the Rings is of course a fundamentally religious and Catholic work; unconsciously so at first, but consciously in the revision. That is why I have not put in, or have cut out, practically all references to anything like `religion', to cults or practices, in the imaginary world. For the religious element is absorbed into the story and the symbolism.
\end{quotex}


Are they horribly misreading LOTR or have they discovered some deep convergence? If there is any interest in this topic, I'll ask for permission to translate the entire post. To me, it seems that Tolkien has affinity with the ancient pagans in the sense that the religious sensibility is integrated into the daily fabric of life; it is not seen as an activity, or belief, separate from one's other daily activities and interests.

\hfill

\texttt{Matt on 2012-04-13 at 01:35 said:}

I'm not in disagreement with you there, was merely acknowledging why they mentioned Hamsun as the worldview in his books seems to be in the same spirit as that of neo-paganism.

But yes, Tolkien is another matter. Eru Illuvatar (The One, Father of All) is a notable example of how the work can not merely be interpreted in a naturalistic sense, as Eru transcends nature and the cosmos.

As to why they would interpret the The Rings books in this way? I don't know what the sites motives are (as I don't even know which site it is), so I can't say for certain, but the reason may be the same as it for many modern minds... .leaving out elements that are significant in a work in order to placate to the pre-conceived opinions of followers.

\hfill

\texttt{Cologero on 2012-04-19 at 00:03 said:}

I have been reluctant to link to \textit{website}\footnote{\url{http://www.corrente88.net/2012/01/17/hamsun-tolkien-lario-e-la-natura/}}, but I will since I just found an \textit{English translation}\footnote{\url{http://oregoncoug.wordpress.com/2012/04/04/hamsun-tolkien-the-aryan-and-nature/}} (allegedly by a \textit{third order Franciscan}\footnote{\url{http://en.wikipedia.org/wiki/Secular_Franciscan_Order}}). The translation is good enough. He misquotes Tolkien, which should be: ``I don't know half of you half as well as I should like; and I like less than half of you half as well as you deserve.'' Also, one phrase makes no sense as translated: it should be ``salvation (\emph{from} what?)''.

I want to point out the profound difference between the consciousness of these neo-pagans compared to that of the pagan Romans as described by Evola in the most recent translation, and to make it a point of discussion.

The Romans did not live in a world of nature, but in a world saturated with occult, divine, and transcendental forces. The Roman experienced a cosmic order and recognized his duty to conform to it. To fail was to commit a \emph{nefas}; salvation was certainly necessary.

This creates a divided consciousness, as every action must relate to the cosmic order, every decision relate to a divine will, every choice relate to one's duties. The true pagan, that is the pagan as a man of Tradition, surpasses this division by transcendence. It seems, though, that the neo-pagan wants to deny the divided consciousness in the first place. He experiences a divine law as a burden, as an impingement on his freedom, so he seeks to live instead in immediacy like the animals. He wants to be part of his environment rather than to struggle against disorder.

A similar narrative, albeit in a intellectually more sophisticated form, is that of the New Right. It, too, denies transcendence, which it can only interpret as totalitarianism. To admit a cosmic or divine law is to admit all must follow it. That is true in one sense. Evola even uses that word to describe the inseparability of the spiritual and the temporal in ancient Rome. All aspects of life are subject to spiritual authority and political power.

Yet, it also leads to perfect freedom, because a man also has his own unique and transcendental dharma to follow, something truly his own. Without a sense of transcendence, what can freedom mean other than to follow one's whims? Rather than a life of freedom, the neo-pagan path leads to a life dominated by material, environmental, and historical factors.

\hfill

\texttt{Matt on 2012-04-19 at 12:55 said:}

Thanks for the link. I read the article on the Italian website (through a translation program on my pc). Its of the usual neo-pagan variety. And with national socialist symbols being placed at the forefront of the site (along with Hitler) I understand your reluctance of posting that link.

As you were giving your summation of the neo-pagan worldview in contrast to the actual pagan worldview of the ancients (or even of a pagan of today who doesn't share the neo-pagan worldview), the connection with gnosticism becomes apparent (which I didn't see before), mostly in terms of the neo-pagan seeing the divine law as a type of totalitarianism. So although there is a love of material nature there that isn't shared with gnostics, there is nonetheless an agreement they have when it comes to the divine law.

\hfill

\texttt{Matt on 2012-04-19 at 12:59 said:}

Thanks for the link. Just slightly browsing that site, and I can see why you were reluctant to post the website. Anyway, the article is of the usual neo-pagan variety.

In your summation of the neo-pagan worldview in contrast to the actual pagan worldview of the ancients, the connection with gnosticism becomes apparent (one that I hadn't quite noticed before), in the view of seeing the divine law as a totalitarianism.

\hfill

\texttt{Charlotte Cowell on 2012-04-20 at 08:35 said:}

I love LOTR. It took me several attempts to get past the first 100 pages then suddenly the book took off and carried me with it, I was unable to put it down, I think it's one of the most amazing things ever written. I have always tended to see it as a Christian work, albeit implied not explict, in the grail tradition -- marriage of the true king with the faerie queen? pretty straightforward from that point of view I would say, no conflict with `pagan' (ie, nature) traditions, just Celtic. A lot of the feeling and description of the Shire comes from Tolkein's dismay at the industrialisation of England following the wars and the attendant war machinery. Really there is no incompatibility whatsoever between the true Christian stream and true religion of former times... ..they are the same

\hfill

\texttt{Cologero on 2012-04-20 at 09:56 said:}

Perhaps more so in the Franciscan tradition? Talking to animals, and so on.

The nun \textit{backs you up}\footnote{\url{http://www.gornahoor.net/?p=1400}}, Charlotte. I think, though, we are more comfortable with seeing a ``development'', not in the sense of evolving and leaving the old behind, but in the sense of a ``deepening'' and an incorporating of earlier elements.

\hfill

\texttt{Charlotte Cowell on 2012-04-20 at 16:22 said:}

Franciscan, yes that's right, he's one of `us' :-) I hadn't really considered whether it is more an evolving or a `deepening', it just seems to be the way it is. As you know I was recently in Central America and I found that contact with the Maya was an extraodinary and thought-provoking experience. now THAT is a strong culture, to have survived so long, throughout countless bloodthirsty revolutions, right up until the present day, spiritually intact. The spirit contacts there are fully manifest on the earth plane and very connected with the earth -- living so closely with nature is such a source of power, it was good to be reminded just how much. It is also interesting to note the degree of syncretism going on, with Baroque catholic churches cheerily sharing space with blatantly pagan forms of worship -- merengue hymns, anyone, pineapple offering and corn shaped into lizards, life-like effigy's of Christ, etc. Also fascinating was how the devoutly Christian element of society has no conflict with the partisan communist guerrillas, who rightly saw in Jesus the blueprint for true revolutionaries... .scary word, but the key is in the `true'. But coming from a Celtic/Druidic/Grail legend perspective, it was very interesting to see so far back into another culture -- to see into their spiritual space, the other side of Atlantis.

\hfill

\texttt{Charlotte Cowell on 2012-04-20 at 16:24 said:}

and yes, I have to agree with the nun -- this is the first thing I learned, this universality -- everything is pure to the pure?

\hfill

\texttt{Cologero on 2012-04-26 at 23:18 said:}

Consider this from Boethius' \textit{Consolation of Philosophy}:

\begin{quotex}\footnotesize

It is well known, and you have seen it demonstrated by astronomers, that beside the extent of the heavens, the circumference of the earth has the size of a point; that is to say, compared with the magnitude of the celestial sphere, it may be thought of as having no extent at all.
\end{quotex}


So another ignorant, modern meme is destroyed, namely, that the ancients believed the earth was the center of the universe and that modern men became disoriented with ``new'' scientific discoveries.

\hfill

\texttt{Charlotte on 2012-04-27 at 05:01 said:}

I am not sure how this meme thread relates to the nun, maybe I'm being a bit dense?

so much of all this is a matter of perspective. From Earth -- where we live, which must count from something in this equation -- the moon and sun appear to be of identical size. Indeed, the moon can fully eclipse the sun at certain times. Yet we `know' (now) that the sun is much larger than the moon and we cannot, therefore, believe our eyes. How strange! They ARE the same size and yet the are NOT! Does not the indisputable evidence seen by those same eyes reveal a great mystery? When spiritually awakened, the human soul is like the moon to the sun of divinity, which shines upon it and is given reflection, by the light of which it must complete its work of atonement. This inner enlightenment habitually leads the imperfect microcosmic moon self to perceive that it is as brilliant as the sun -- the `same size', so to speak -- by simple virtue of the fact that it shines where previously it was darkened. Of course, a wise soul `knows', intrinsically, that for all these fantastic perceptions, which give all the appearances and sensation of highest reality, it CANNOT in fact be of comparable brilliance to the divine light shining upon it. Logically this is impossible. The soul can only do its best to make the mirror of mysteries as clean as possible, thereby ensuring a true reflection when the Holy Spirit turns to gaze upon it. In this situation the soul cannot `believe its eyes', it has to trust to a deeper wisdom knowledge that speaks truth directly to the heart, its core self and uncorrupted spark. The signs in the skies are there to help us understand the nature and pathways of our souls -- and very simple to understand if one doesn't think too hard! As for Earth being centre of the universe... ..well in a certain sense it IS, if one looks `down and in' upon the cosoms, rather than `up and out' from here. Something to do with the elipses and angles taken by the various heavenly bodies,all moving in arcs and triangles into strange places and funny directions for all kinds of distances. Earth from the higher vantage points is like a glowing pearl around which everything revolves, like a precious egg. The first astronauts marvelled at how Earth appeared from space, remarking that in contrast with all the other planets it was shining brightly with a brilliant aura, the light of the world soul, apple of God's eye.

\hfill

\texttt{Angolmois on 2012-04-28 at 10:38 said:}

I think it was in RUNA in which there was an article -- perhaps by Thorsson or Moynihan -- that linked LOTR and Tolkien with `Radical Traditionalism', and I think this is quite a justified assumption. Yet, I do agree myself that there is a convergence of Pagan and Christian wisdom \& symbolism of the highest degree in LOTR. The transference of `Gandalf the Grey' into `Gandalf the White' can be seen as an aeonic transference from Heathenism to Christianity of the Traditional kind. And may I remember that it was the blasphemed outcast who became `The Golden-red King' and the unifier of different races / traditions. Tolkien was a genius, one of a kind. His Catholicism reeks Paganism more than Dante's variety does.

The earth IS the centre of the universe, but only as a stepping stone and as a foothold!

\hfill

\texttt{Charlotte on 2012-04-28 at 11:16 said:}

Interesting, we seem to agree on a lot to do with LOTR!

Regarding characters, tho, I had tended to see the King as being Christlike (the faerie bride his Holy Soul), with Gandalf in a more Holy Fatherly role? However, his transformation from grey to white is indeed telling, as you say. I wonder if this transformation is supposed to give a clue to the initiate about the changing essence of their interactions with the `force', the divine One? In any event, the passing away of the old world is something we all have to face, alogn with the sacrifice of so many Shining Ones. Thank heaven for the sacred Earth given Sam-wise to tend, not to mention the benevolent ruling messiah, complete with bride... ..

\hfill

\texttt{Angolmois on 2012-04-28 at 13:25 said:}

Of course the King is Christ; he wears the purple-red-robe and a golden crown, the symbols of royalty. Gandalf as a high priest relates to the pontificate, and thus LOTR is in agreement with Guénon rather than Evola. Internally we (must) have both Gandalf and Aragorn, together with others, within us in harmony and symbiosis. Saruman is clearly the fallen Lucifer, downtrodden by mans hubris, the black tower of pride and the main centre of mechanical titanic force, and Sauron is of course no other than Satan himself, da.

As regards ``geocentrism'' via a heretic's wormhole, this must be one of my favourite scenes from the LOTR epic:

\url{http://www.youtube.com/watch?v=6cdc_OaAzQM&feature=fvwrel}


Balrog ``the flame of Udun'' reminds me of The Primordial Fire UR (the flame of Ur could be his name), ``that old devil'', a primordial material and animalistic force personified, whereas the brigde reminds me of the Antahkarana-rainbowbrigde in which the old ego must die in order to be born again. Gandalf's fight in the crumbling brigde together with the flaming sword and the magicians staff clearly symbolise the two-fold functions of Pater Jupiter and Venus-Lucifer, only the sword and staff are in wrong hands, but that's not important. Gandalf's fight with Balrog is his final battle with his own shadow and his lower animalistic self; he no longer has any human vices when he again appears in the scenery as Gandalf the White. And so on.

\hfill

\texttt{Charlotte on 2012-04-28 at 13:55 said:}

ah I see what you mean now, makes me wonder what this symbol means for the papacy as it is now... .the final pope is the one to take up the white robe, then? Part of me still thinks JPII went onto fulfill this function, another part of me waits in pensive doubt. Church is being purified from within but they can't withold answers forever. If they fail to understand themselves, they should say.

One thing I wondered about the film is why they missed out Tom Bombadil, which IMO was one of the best characters/episodes in the early book. I guess time ran out, but I feel it's a section that needs revisting, the early part is very dense.

My own favourite scene, I think, is when Arwen calls up the river horses to keep the black riders at bay... good invocatory skills, perfect timing, great example :-)

\hfill

\texttt{Angolmois on 2012-04-28 at 13:58 said:}

It seems to me that what you fear is a cage, Charlotte.

\hfill

\texttt{Charlotte on 2012-04-28 at 14:39 said:}

well you are right there, everytime I visit I zoo that strikes me forcibly. Interesting to learn while I was in Guatemala that their national bird, the fabled, beautiful Quetzal, is the one creature that simply cannot survive in captivity -- it dies when caged, which is fitting for something that signifies liberation above all else. Something about the truth, a little bird said, that is also meant to set one free.

\hfill

\texttt{Matt on 2012-04-28 at 19:04 said:}

Angolmois,

Its been awhile since I've read the books and The Silmarillion, but isn't Melkor supposed to be Lucifer? Though I suppose once could say that all three figures (Melkor, Sauron, and Saruman) represent the same archetype.

\hfill

\texttt{Matt on 2012-04-28 at 19:07 said:}

Oh and sorry for my earlier somewhat double post. Just noticed it in now. The first post didn't seem to register (but so I posted a shortened version.

\hfill

\texttt{Angolmois on 2012-04-29 at 03:12 said:}

I haven't read Silmarillion myself (shame on me!) but I agree with you that there are many representatives of the same luciferian archetype in Tolkien's books. The monster Balrog can be traced back to this ``dark fire'' (Boehme) also.

\hfill

\texttt{strangelove on 2012-05-15 at 15:15 said:}

More info on our geocentric universe here:

\url{http://christian-wilderness.forumvi.com/t39-stationary-earth}


\end{sffamily}
\end{footnotesize}

